\documentclass[11pt,a4paper]{article}
\usepackage[T1]{fontenc}
\usepackage{lmodern}
\usepackage[margin=25mm]{geometry}
\usepackage{amsmath,amssymb,amsthm,mathtools}
\usepackage[hidelinks]{hyperref}
\usepackage{microtype}
\newcommand{\Q}{\mathbb Q}
\newcommand{\R}{\mathbb R}
\newcommand{\N}{\mathbb N}
\title{An analytic $5$-adic inverse with arbitrarily small nonconvergent orbits}
\author{Submission for TLMC problem 00000002704}
\date{24 October 2026}
\begin{document}
\maketitle

\paragraph{Claim and scope.}
The conjecture says that iteration of a $p$-adic local inverse within its
convergence radius converges linearly. No attracting multiplier hypothesis
is stated. We give an actual one-dimensional analytic inverse over $\Q_5$
whose power series has infinite convergence radius, but which has
nonconvergent orbits entirely inside every positive-radius ball around its
fixed point. This refutes the unconditional convergence clause. It does not
contradict the inverse function theorem or a contraction theorem with its
attracting hypotheses, and does not assign a value to an unspecified
linear rate or logarithmic constant.

\paragraph{The inverse and the actual power series.}
Let
\[
f:\Q_5\longrightarrow\Q_5,\quad f(z)=-z,
\qquad g:\Q_5\longrightarrow\Q_5,\quad g(z)=-z.
\]
Both maps are continuous and $\Q_5$-linear. Direct substitution gives
\[
g\circ f=f\circ g=\operatorname{id}_{\Q_5}.
\]
They therefore define a global homeomorphism and its actual inverse;
restriction to any invariant neighborhood is a local inverse pair.
The derivative is
\[
f'(z)=-1\ne0,\qquad |f'(z)|_5=1
\]
at every point, so the one-dimensional Jacobian is invertible.

The inverse has the genuine series at $0$
\[
g(z)=0+(-1)z+0z^2+0z^3+\cdots.
\]
This is a polynomial series, hence has convergence radius $+\infty$.
The formal proof uses the native
\texttt{ContinuousLinearMap.fpowerSeries} of the actual inverse map and
proves its native \texttt{FormalMultilinearSeries.radius} is $\infty$.
It also supplies \texttt{HasFPowerSeriesOnBall} on that whole ball, so the
series is proved to represent $g$, rather than merely being a formal list of
coefficients. These are the linear analytic constructions of Mathlib
\cite{linear}.

\paragraph{Nonzero starting points in arbitrarily small balls.}
Fix any real $r>0$. The normalized $5$-adic norm satisfies
\[
|5|_5=\frac15,\qquad |5^N|_5=5^{-N}.
\]
Choose $N\in\N$ with $5^{-N}<r$ and put $x=5^N$.
Because $\Q_5$ is a field of characteristic zero, $x\ne0$, and
\[
x\in B(0,r).
\]
All these quantities come from the actual $5$-adic field and its norm;
the point and its size are not assigned abstractly.

\paragraph{The entire orbit stays in the ball and fails to converge.}
For every $n\in\N$, direct iteration gives
\[
g^{2n}(x)=x,\qquad g^{2n+1}(x)=-x,
\qquad |g^n(x)|_5=|x|_5<r.
\]
Thus the entire orbit remains in $B(0,r)$, and in particular remains within
the actual infinite convergence radius of the inverse series. This holds
for every $r>0$, not just for one large neighborhood.

Suppose the full orbit converged to some $y\in\Q_5$. The maps
$n\mapsto2n$ and $n\mapsto2n+1$ both tend to infinity, so the even and odd
subsequences would converge to the same $y$. They are respectively the
constant sequences $x$ and $-x$. Uniqueness of limits in the Hausdorff field
$\Q_5$ would give
\[
x=y=-x,
\qquad 2x=0.
\]
Since $2\ne0$ and $x\ne0$ in $\Q_5$, this is impossible. Therefore the orbit
does not converge at all, and consequently cannot converge linearly.

\paragraph{Formal scope and files.}
The Lean conclusion is quantified over every real $r>0$. It constructs
an actual nonzero $x\in\Q_5$, proves every iterate lies in the open metric
ball $B(0,r)$, and excludes convergence to every $y\in\Q_5$. Separate
theorems prove both inverse identities, the global homeomorphism, the strict
derivative, its nonvanishing and norm, the actual series representation and
its infinite radius, and the two cofinal constant subsequences.

The complete construction is in \texttt{lean/Main.lean}. Its conclusion is
\begin{center}
\texttt{TLMC2704.counterexample}.
\end{center}
The project pins Lean 4.33.0 and Mathlib commit
\begin{center}
\small\texttt{db584cd6d46c92f209a44c0f1c829460d327499d}.
\end{center}
The original bilingual conjecture is preserved in \texttt{SOURCE.md}.
Exact build commands, exhaustive namespace axiom checks, fresh-kernel
receipts and file hashes accompany the submission. The proof addresses
the stated convergence assertion only; it makes no claim that multiplier
norm alone determines a general analytic inverse's radius.

\begin{thebibliography}{1}
\bibitem{linear}
Mathlib, \emph{Linear functions are analytic},
\texttt{Mathlib/Analysis/Analytic/Linear.lean},
\href{https://github.com/leanprover-community/mathlib4/blob/db584cd6d46c92f209a44c0f1c829460d327499d/Mathlib/Analysis/Analytic/Linear.lean}{pinned source}.
\end{thebibliography}
\end{document}
